\chapter{HASIL DAN ANALISIS}

% TODO: Isi bab ini setelah eksperimen (jadwal: September--Oktober 2026) selesai
% dan data dianalisis. Setiap subbab di bawah merupakan kerangka yang
% disusun berdasarkan desain eksperimen pada Bab III.

\section{Deskripsi Lingkungan Eksperimen}

% TODO: Dokumentasikan kondisi aktual saat eksperimen berjalan:
% - Spesifikasi klaster (3 node, 24 inti CPU, ~23 Gi memori, Kubernetes v1.29+)
% - Versi ArgoCD, interval sync (60 detik), self-heal dan prune aktif
% - Beban uji workload komposit multi-service (>= 16 komponen)
% - Profil operator dan hasil practice run
% - Variabel pengganggu yang terpantau (CPU < 70%, RTT, dsb.)

\section{Hasil E1: Konsistensi \textit{Deployment} (R1)}

% TODO: Sajikan hasil pengukuran persentase sumber daya yang cocok dengan
% desired declarative state pada t0, t1, dan t2 untuk setiap skenario (A--G)
% pada kedua lingkungan. Gunakan tabel dan/atau box plot per skenario.

\section{Hasil E2: Waktu Deteksi, Pemulihan, dan Identifikasi Akar Masalah (R2)}

% TODO: Sajikan hasil pengukuran waktu deteksi (td - t0), waktu pemulihan
% (t2 - t0), dan waktu identifikasi akar masalah (trci - td) untuk setiap
% skenario pada kedua lingkungan (n = 10 per skenario). Sertakan juga
% jumlah perintah kubectl pada Lingkungan A sebagai pengamatan pelengkap.

\section{Hasil E3: \textit{Traceability} Operasional (R3)}

% TODO: Dokumentasikan secara deskriptif kelengkapan audit log untuk kedua
% lingkungan: Git commit history + ArgoCD notifications (Lingkungan B) versus
% Kubernetes audit logs (Lingkungan A), serta kemudahan identifikasi akar
% masalah dan keterkaitannya dengan R2.

\section{Statistik Deskriptif}

% TODO: Sajikan mean, median, rentang, dan deviasi standar untuk setiap
% metrik per lingkungan per skenario, dalam bentuk tabel ringkasan.

\section{Uji Perbandingan \textit{Mann--Whitney U}}

% TODO: Sajikan hasil uji Mann--Whitney U untuk setiap metrik antara
% Lingkungan A dan Lingkungan B pada tingkat signifikansi alpha = 0,05.
% Sertakan nilai p untuk setiap skenario.

\section{Ukuran Efek}

% TODO: Sajikan nilai Cohen's d (atau Cliff's delta) untuk setiap metrik dan
% interpretasikan mengikuti konvensi Cohen: |d| < 0,2 (kecil), ~0,5 (sedang),
% > 0,8 (besar). Identifikasi skenario dengan efek terbesar dan terkecil.

\section{Interpretasi Hasil}

% TODO: Interpretasikan hasil statistik dalam konteks literatur (Bab II) dan
% praktik operasional. Jawab dimensi evaluasi pada Bab I: konsistensi,
% waktu deteksi/pemulihan/identifikasi akar masalah, dan traceability.

\section{Pembahasan}

% TODO: Bahas skenario penyimpangan yang paling berdampak, faktor yang
% memengaruhi hasil, kondisi di mana manual deployment tetap memadai, dan
% implikasi praktis bagi pemilihan paradigma deployment. Kaitkan dengan
% ancaman validitas yang diidentifikasi pada Bab III.
